\cotefile{77}
\scanpage{18}
Quant aux formes quadratiques invariantes \inserted{quasi-inv.} sur
l'espace des translations $\check V$ (toujours sous l'hyp.
$\alpha^2-4$ inv) elles sont comprises parmi les
\unsure{quasi-invar.|quasi-inva.|quasi-inv.} sur \unsure{$\widetilde E$|$t^{=}$|$\check E$} (dont \inserted{\deleted{\unsure{l'invariance pour}}\deleted{\unsure{une certaine}} puisque $\alpha\ne2$ fibre par fibre est un} quotient $\unsure{\}|)}$,
on trouve donc, si $\alpha$ inv., les multiples
seulement de la forme $xy$ (qui est non
dégénérée et donne lieu à la conique circonscrite
à $P$), et si $\alpha=0$ i.e. \deleted{\unsure{que}}
${\zeta}^2=-1$ \inserted{(donc $2$ inv. sur $S$)}, on trouve comme dessus \deleted{les}
multiples de $x^2+y^2$ et de $x^2-y^2$, dont
la \deleted{\unsure{deuxième devant prendre les mêmes valeurs}}\inserted{première prend les valeurs}
$1+1=2$ et $\zeta^2+{\zeta'}^2=\alpha^2-2=-2$ sur les sommets
de $P$ \inserted{(donc deux valeurs distinctes)}, et la deuxième y prend la valeur $0$,
mais (comme nous avons remarqué) définit une
conique circonscrite dégénérée\ldots .

\begin{center}\rule{1.6cm}{0.3pt}\end{center}

Il reste à étudier la question des formes
quadratiques \inserted{quasi-}invariantes, d'abord pour $\alpha=2$
et $\alpha=-2$, puis éventuellement pour $\alpha$
``variable'' général\unsure{?}\ldots .

% archive sheet number: 17

